\exercise{7.4}{Rational and irrational terms can have different kinds of limits}
\goal Give \textup{(a)} a sequence of irrational numbers converging to a
rational number, and \textup{(b)} a sequence of rational numbers
converging to an irrational number.

\idea For part \textup{(a)}, divide $\sqrt2$ by larger positive integers.
For part \textup{(b)}, use the lower decimal truncations of $\sqrt2$.
We verify both the nature of every term and convergence.

\step{1. Recall why $\sqrt2$ is irrational}
The irrationality of $\sqrt2$ was established in Section~2.
Here is the short integer argument as well. If
$\sqrt2=p/q$ in lowest terms, with integers $p,q$ and $q>0$,
then $p^2=2q^2$, so $p^2$ is even. The square of an odd integer
$2j+1$ is $4j(j+1)+1$, which is odd. Thus $p$ must be even;
write $p=2j$. Substitution gives
\[
 4j^2=2q^2,\qquad q^2=2j^2.
\]
Now $q^2$ is even, so the same odd-square argument makes $q$ even.
The numerator and denominator would have the common factor $2$,
contrary to lowest terms. Hence $\sqrt2\notin\Q$.

\step{2. Part (a): Irrational terms with rational limit}
Define, for every $n\in\N$,
\[
 x_n=\frac{\sqrt2}{n}.
\]
Since $n\ge1$, division is legitimate. If any term $x_n$ were
rational, multiplication by the rational integer $n$ would make
$nx_n=\sqrt2$ rational. This contradicts Step~1.
Thus every term $x_n$ is irrational.

To prove convergence, let $\varepsilon>0$ and choose an integer
$N>\sqrt2/\varepsilon$ by the Archimedean property. For $n>N$,
\[
 \abs{x_n-0}=\frac{\sqrt2}{n}<\varepsilon.
\]
Therefore $x_n\to0$. The limit is rational because $0=0/1$.

\step{3. Part (b): Define each decimal truncation using a finite maximum}
For a real $t\ge0$, the notation $\lfloor t\rfloor$ denotes the
integer $m$ satisfying
\[
 m\le t<m+1.
\]
To justify existence, choose a positive integer $M>t$ using the
Archimedean property. The finite set
\[
 F=\{j\in\Z:0\le j\le M\text{ and }j\le t\}
\]
is nonempty because it contains $0$. Let $m=\max F$, which exists
for a nonempty finite set of integers. Then $m\le t$.
If $m+1\le t$, the inequalities
$0\le m+1\le t<M$ would give $m+1\in F$, contradicting the
maximality of $m$. Hence $t<m+1$, as required.

This integer is unique. If two integers $m<j$ satisfied their
respective inequalities, then $m+1\le j\le t$ would contradict
$t<m+1$. Interchanging them handles the other possible order.
Thus the floor notation is well-defined here.

Apply this construction to $t=10^n\sqrt2$ and put
\[
 m_n=\left\lfloor10^n\sqrt2\right\rfloor,
 \qquad r_n=\frac{m_n}{10^n}\quad(n\in\N).
\]
Every $m_n$ is an integer and $10^n$ is a positive integer.
Therefore every $r_n$ is rational by the definition of $\Q$.
The defining floor inequalities give
\[
 m_n\le10^n\sqrt2<m_n+1.
\]
Dividing by $10^n>0$ preserves order and yields
\[
 r_n\le\sqrt2<r_n+\frac1{10^n}.
\]
After subtracting $r_n$, this becomes the precise truncation error
estimate
\[
 0\le\sqrt2-r_n<\frac1{10^n}.
\]
In fact the left inequality is strict, because a rational $r_n$
cannot equal the irrational number $\sqrt2$.

\step{4. Prove convergence of the rational truncations}
We may compare the error with $1/n$. The inequality $10^n\ge n$
holds for all positive integers $n$: the base case is $10\ge1$;
if it holds for $n$, then
\[
 10^{n+1}=10\cdot10^n\ge10n\ge n+1,
\]
because $10n-(n+1)=9n-1\ge8>0$ for $n\ge1$.
Induction proves the inequality. Taking reciprocals of positive
quantities gives $1/10^n\le1/n$.

Let $\varepsilon>0$ and choose an integer $N>1/\varepsilon$.
For $n>N$, we obtain
\[
 \abs{r_n-\sqrt2}
 =\sqrt2-r_n<\frac1{10^n}\le\frac1n<\varepsilon.
\]
Thus $r_n\to\sqrt2$, and its limit is irrational by Step~1.

\answer{For \textup{(a)}, take $x_n=\sqrt2/n\to0$, with every
term irrational and the limit rational. For \textup{(b)}, take
$r_n=\lfloor10^n\sqrt2\rfloor/10^n\to\sqrt2$, with every term
rational and the limit irrational.}

\exercise{7.5}{Three differences of large square roots}
\goal Find the limits of
\[
 \text{(a)}\ \sqrt{n^2+1}-n,\qquad
 \text{(b)}\ \sqrt{n^2+n}-n,\qquad
 \text{(c)}\ \sqrt{4n^2+n}-2n.
\]

\idea Multiply each difference by its conjugate divided by that
same conjugate. The difference-of-squares identity removes the
subtraction in the numerator. We then bound the error from the
proposed limit directly, so no later limit laws are needed.

\step{1. Part (a): Expand the product with the conjugate}
Let $a_n=\sqrt{n^2+1}-n$. For $n\ge1$, the square root is
nonnegative and $n>0$, so the conjugate
$\sqrt{n^2+1}+n$ is positive. Dividing by it is valid.
Using $(u-v)(u+v)=u^2-v^2$ gives
\begin{align*}
 a_n
 &=\frac{(\sqrt{n^2+1}-n)(\sqrt{n^2+1}+n)}
         {\sqrt{n^2+1}+n}\\
 &=\frac{(\sqrt{n^2+1})^2-n^2}{\sqrt{n^2+1}+n}\\
 &=\frac{n^2+1-n^2}{\sqrt{n^2+1}+n}\\
 &=\frac1{\sqrt{n^2+1}+n}.
\end{align*}
Since $n^2+1>n^2$ and $n\ge0$, the nonnegative square root
$\sqrt{n^2+1}$ is greater than $n$. Indeed, the contrary inequality
would give $n^2+1\le n^2$ upon squaring nonnegative quantities.
Therefore
\[
 0<a_n=\frac1{\sqrt{n^2+1}+n}\le\frac1{2n}.
\]
Given $\varepsilon>0$, choose an integer $N>1/(2\varepsilon)$.
For $n>N$, this estimate gives
\[
 \abs{a_n-0}=a_n\le\frac1{2n}<\varepsilon.
\]
Thus the limit in part \textup{(a)} is $0$.

\step{2. Part (b): Rationalize and divide by the positive index}
Write $b_n=\sqrt{n^2+n}-n$. Its conjugate is positive for
$n\ge1$, and the same multiplication gives
\begin{align*}
 b_n
 &=\frac{(\sqrt{n^2+n}-n)(\sqrt{n^2+n}+n)}
         {\sqrt{n^2+n}+n}\\
 &=\frac{(n^2+n)-n^2}{\sqrt{n^2+n}+n}\\
 &=\frac n{\sqrt{n^2+n}+n}.
\end{align*}
The radicand factors as $n^2+n=n^2(1+1/n)$. Because $n>0$,
\[
 \sqrt{n^2+n}=n\sqrt{1+\frac1n}.
\]
To check this identity, the right-hand side is nonnegative, and
its square is $n^2(1+1/n)=n^2+n$; uniqueness of the nonnegative
square root then gives the equality. Define
$q_n=\sqrt{1+1/n}$. Substitution and cancellation of $n\ne0$ yield
\[
 b_n=\frac n{nq_n+n}
     =\frac n{n(q_n+1)}
     =\frac1{q_n+1}.
\]
Here $q_n>1$, since $q_n\ge0$ and
$q_n^2=1+1/n>1$; a nonnegative $q_n\le1$ would have square
at most $1$. Hence $q_n+1>2$, so $0<b_n<1/2$.

\step{3. Part (b): Compute the exact error from $1/2$}
All the denominators below are positive. Since $b_n<1/2$,
\begin{align*}
 \abs{b_n-\frac12}
 &=\frac12-\frac1{q_n+1}\\
 &=\frac{q_n+1-2}{2(q_n+1)}\\
 &=\frac{q_n-1}{2(q_n+1)}\\
 &=\frac{(q_n-1)(q_n+1)}{2(q_n+1)^2}\\
 &=\frac{q_n^2-1}{2(q_n+1)^2}\\
 &=\frac{(1+1/n)-1}{2(q_n+1)^2}\\
 &=\frac1{2n(q_n+1)^2}
 \le\frac1{8n}.
\end{align*}
The last inequality uses $(q_n+1)^2\ge2^2=4$ and $n>0$.
For any $\varepsilon>0$, choose an integer $N>1/(8\varepsilon)$.
Then $n>N$ implies
$\abs{b_n-1/2}\le1/(8n)<\varepsilon$.
Therefore the limit in part \textup{(b)} is $1/2$.

\step{4. Part (c): Rationalize and factor the square root}
Put $c_n=\sqrt{4n^2+n}-2n$. Since $n\ge1$, its conjugate
$\sqrt{4n^2+n}+2n$ is positive. Consequently
\begin{align*}
 c_n
 &=\frac{(\sqrt{4n^2+n}-2n)(\sqrt{4n^2+n}+2n)}
         {\sqrt{4n^2+n}+2n}\\
 &=\frac{(4n^2+n)-(2n)^2}{\sqrt{4n^2+n}+2n}\\
 &=\frac{4n^2+n-4n^2}{\sqrt{4n^2+n}+2n}\\
 &=\frac n{\sqrt{4n^2+n}+2n}.
\end{align*}
Now $4n^2+n=n^2(4+1/n)$. The positive-index square-root
identity gives
\[
 \sqrt{4n^2+n}=n\sqrt{4+\frac1n},
\]
because both sides are nonnegative and the square of the right-hand
side is $n^2(4+1/n)=4n^2+n$. Let
$p_n=\sqrt{4+1/n}$. Canceling the nonzero factor $n$ gives
\[
 c_n=\frac n{np_n+2n}
     =\frac n{n(p_n+2)}
     =\frac1{p_n+2}.
\]
Since $p_n\ge0$ and $p_n^2=4+1/n>4$, we have $p_n>2$.
Thus $p_n+2>4$ and $0<c_n<1/4$.

\step{5. Part (c): Compute the exact error from $1/4$}
Using $c_n<1/4$ and then multiplying numerator and denominator
by the positive number $p_n+2$, we obtain
\begin{align*}
 \abs{c_n-\frac14}
 &=\frac14-\frac1{p_n+2}\\
 &=\frac{p_n+2-4}{4(p_n+2)}\\
 &=\frac{p_n-2}{4(p_n+2)}\\
 &=\frac{(p_n-2)(p_n+2)}{4(p_n+2)^2}\\
 &=\frac{p_n^2-4}{4(p_n+2)^2}\\
 &=\frac{(4+1/n)-4}{4(p_n+2)^2}\\
 &=\frac1{4n(p_n+2)^2}
 \le\frac1{64n}.
\end{align*}
Here $(p_n+2)^2\ge4^2=16$, so the denominator is at least
$4n\cdot16=64n$. For any $\varepsilon>0$, choose an integer
$N>1/(64\varepsilon)$. Then $n>N$ gives
$\abs{c_n-1/4}\le1/(64n)<\varepsilon$.
This proves that the limit in part \textup{(c)} is $1/4$.

\answer{The limits are \textup{(a)} $0$, \textup{(b)} $1/2$,
and \textup{(c)} $1/4$. The displayed conjugate identities and
error estimates justify all three values directly.}
